\section*{Chapitre \ref{ch:b3:cell-cycle-apoptosis} --- Contrôle du cycle cellulaire et mort cellulaire programmée}

\begin{solution}{exo:b3:cell-cycle-apoptosis:1}
G1, S, G2, M. Progression en G1 et \omterm{def:b3:cell-cycle-apoptosis:checkpoints}{point de restriction}~: \omterm{def:b3:cell-cycle-apoptosis:cdk}{cycline}
D--Cdk4/6~; G1/S~: \omterm{def:b3:cell-cycle-apoptosis:cdk}{cycline} E--Cdk2~; S~: \omterm{def:b3:cell-cycle-apoptosis:cdk}{cycline} A--Cdk2~; G2/M~:
\omterm{def:b3:cell-cycle-apoptosis:cdk}{cycline} B--Cdk1. Le \omterm{def:b3:cell-cycle-apoptosis:cdk}{complexe promoteur de l'anaphase} (APC/C) détruit la
\omterm{def:b3:cell-cycle-apoptosis:cdk}{cycline}~B et la sécurine et met fin à la mitose.
\end{solution}

\begin{solution}{exo:b3:cell-cycle-apoptosis:2}
Hartwell~: les mutants \emph{cdc} de la levure bourgeonnante, arrêtés à
des stades définis, et la notion de Start. Nurse~: \emph{cdc2} de la
levure à fission comme kinase qui fixe le moment de la mitose, et son
homologue humain (Cdk1) qui sauve la levure --- la conservation.
Hunt~: la \omterm{def:b3:cell-cycle-apoptosis:cdk}{cycline} dans les œufs d'oursin, protéine synthétisée tout au
long du cycle et détruite à chaque division.
\end{solution}

\begin{solution}{exo:b3:cell-cycle-apoptosis:3}
\omterm{def:b3:cell-cycle-apoptosis:checkpoints}{Point de restriction} (fin de G1)~: facteurs de croissance, nutriments,
taille~; agit sur \omterm{def:b3:cell-cycle-apoptosis:cdk}{cycline}~D--Cdk4/6 et Rb. G2/M~: ADN intact et
répliqué~; \omterm{def:b3:dna-repair:ddr}{ATM}/ATR inhibent Cdc25, ce qui maintient Cdk1 phosphorylée
et inactive. Assemblage du fuseau~: chaque kinétochore attaché~; les
kinétochores non attachés forment des complexes Mad2--Cdc20 qui
inhibent l'APC/C.
\end{solution}

\begin{solution}{exo:b3:cell-cycle-apoptosis:4}
\omterm{def:b3:cell-cycle-apoptosis:apoptosis}{Apoptose}~: la cellule se rétracte, la \omterm{def:b3:chromatin-epigenetics:nucleosome}{chromatine} se condense, l'ADN est
coupé en échelle, la membrane bourgeonne mais reste close, la
\omterm{def:b3:cell-cycle-apoptosis:apoptosis}{phosphatidylsérine} est exposée, le cadavre est englobé, pas
d'inflammation. \omterm{def:b3:cell-cycle-apoptosis:apoptosis}{Nécrose}~: la cellule gonfle, la membrane se rompt, le
contenu fuit, l'ADN est dégradé au hasard, inflammation. Exécutrices~:
les \omterm{def:b3:cell-cycle-apoptosis:apoptosis}{caspases}, protéases à cystéine qui coupent après un aspartate.
\end{solution}

\begin{solution}{exo:b3:cell-cycle-apoptosis:5}
Montée de $5$ à $30$ à $1$ par minute~: \qty{25}{min}. Dégradation de
$30$ à $5$ avec une demi-vie de \qty{2}{min}~: $\log_{2}(30/5)\times 2 =
\qty{5.2}{min}$. Période d'environ \qty{30}{min}, la mitose en faisant
\qty{17}{\%}.
\end{solution}

\begin{solution}{exo:b3:cell-cycle-apoptosis:6}
(a) $C$ ne peut jamais tomber sous $C_{1}$~: bloquée en mitose sur la
branche haute. (b) Pas de phosphorylation inhibitrice~: le seuil
$C_{2}$ est abaissé, la mitose commence tôt à petite taille (le
phénotype «~wee~»). (c) Cdk1 reste phosphorylée~: la branche haute est
hors d'atteinte, arrêt en G2 avec des cellules allongées. (d) Comme en
(b), et le \omterm{def:b3:cell-cycle-apoptosis:checkpoints}{point de contrôle} G2/M, qui agit par Wee1/Cdc25, ne peut
plus retenir la cellule.
\end{solution}

\begin{solution}{exo:b3:cell-cycle-apoptosis:7}
Noyau G1 dans un cytoplasme de phase~S~: ses origines sont autorisées
et le cytoplasme fournit des couples \omterm{def:b3:cell-cycle-apoptosis:cdk}{cycline}~E/A--Cdk2 actifs, donc il
entre aussitôt en S. Noyau G2 dans un cytoplasme de phase~S~: ses
origines ont démarré et n'ont pas été réautorisées (les facteurs
d'autorisation sont détruits ou inhibés par les CDK), donc il ne peut
pas se répliquer de nouveau~; il attend que le partenaire atteigne G2
et tous deux entrent en mitose ensemble.
\end{solution}

\begin{solution}{exo:b3:cell-cycle-apoptosis:8}
Le kinétochore non attaché est un catalyseur~: il convertit sans cesse
Mad2 dans sa conformation active, produisant un inhibiteur diffusible
de Cdc20 qui se répand dans la cellule et maintient l'APC/C éteint
partout --- un seul kinétochore y suffit. Sans Mad2, l'APC/C est
activé dès que Cdk1 l'a allumé, que les chromosomes soient attachés ou
non~: l'anaphase commence avec des chromosomes non attachés, les filles
sont aneuploïdes, et l'organisme (ou la lignée cellulaire) meurt de
pertes de chromosomes.
\end{solution}

\begin{solution}{exo:b3:cell-cycle-apoptosis:9}
Sans \emph{ced-9}~: des cellules qui devraient vivre meurent --- une
létalité embryonnaire par excès de mort. Sans \emph{ced-3}~: aucune des
$131$ morts n'a lieu, les cellules survivent et se différencient, le
ver est viable. Sans les deux~: pas de mort --- le phénotype
\emph{ced-3}. Puisque retirer l'exécuteur annule l'effet du retrait du
gardien, CED-9 agit en amont, en maintenant CED-4/CED-3 inactifs~;
quand il disparaît ils sont actifs, sauf s'ils sont absents eux aussi.
\end{solution}

\begin{solution}{exo:b3:cell-cycle-apoptosis:10}
Avec $A = aC$, $\mathrm{d}C/\mathrm{d}t = k_{s} - k_{d}aC^{2}$, équation
à une variable d'état stationnaire $C^{*} = \sqrt{k_{s}/(k_{d}a)}$ et de
pente $-2k_{d}aC^{*} < 0$ en ce point~: tout écart décroît de façon
monotone, et une seule variable du premier ordre ne peut pas osciller.
La courbe en S donne deux branches stables séparées par une zone
interdite, de sorte que l'état doit sauter et ne peut pas se poser~; un
délai explicite dans la rétroaction donnerait des oscillations par une
autre voie --- le frein arrivant trop tard, le dépassement, et ainsi de
suite --- comme dans bien des rythmes hormonaux.
\end{solution}

\begin{solution}{exo:b3:cell-cycle-apoptosis:11}
Chaque capteur qui arrive est capturé par un gardien, de sorte que rien
ne se passe tant que les $N$ gardiens ne sont pas tous occupés, vers
$t \approx N/r$~; les capteurs suivants sont libres, activent Bax/Bak,
dont la formation de pores croît fortement avec leur nombre, et les
mitochondries s'ouvrent en quelques minutes. Un stimulus plus fort ne
raccourcit le délai que par $r$~; au-delà du seuil l'issue est la même.
Le vénétoclax occupe des gardiens, ce qui réduit le $N$ efficace~: le
délai se raccourcit, et une cellule amorcée --- déjà proche de $N$ ---
meurt aussitôt.
\end{solution}

\begin{solution}{exo:b3:cell-cycle-apoptosis:12}
L'excès de Bcl-2 bloque la \omterm{def:b3:cell-cycle-apoptosis:pathways}{voie intrinsèque}~: des lymphocytes~B qui
devraient mourir à la fin de leurs signaux de croissance survivent, et
la population croît par défaut de mort, à la vitesse lente à laquelle
de telles cellules sont produites --- indolente, jusqu'à ce qu'une
seconde lésion y ajoute de la prolifération. La perte de Rb supprime le
\omterm{def:b3:cell-cycle-apoptosis:checkpoints}{point de restriction}~: les précurseurs rétiniens traversent le cycle
sans facteurs de croissance et se divisent aussi vite que le moteur le
permet --- agressive. Les deux tumeurs sont les deux programmes en
défaut chacun de son côté~: un contrôle de la mort et un contrôle de la
division, dont la perte suffit à faire une tumeur, la seconde plus
vite.
\end{solution}

\begin{solution}{pb:b3:cell-cycle-apoptosis:1}
\textbf{1.} $10^{7}\times 3500/5 = 7\times 10^{9}$ cellules par jour,
environ \num{80000} par seconde.
\textbf{2.} $3500/5 = 700$ par crypte et par jour~; $22\times 2^{5} = 704$.
\textbf{3.} $80\times 365 \approx \num{29000}$ divisions~; $\times 0.6
\approx \num{17500}$ mutations.
\textbf{4.} Les cellules de transit sont éliminées en quelques jours et
emportent leurs mutations~; la cellule souche persiste toute la vie, de
sorte que chacune de ses mutations est gardée et que chaque division en
ajoute --- un cycle lent les minimise.
\textbf{5.} La villosité se vide en à peu près les $5$ jours de son
transit normal~; la rebâtir demande un jour de récupération, cinq
divisions de \qty{12}{h} et la migration le long de la villosité ---
quatre à cinq jours, pendant lesquels la barrière est nue~: les
divisions de transit rapides de l'intestin en font une victime précoce
du rayonnement.
\textbf{6.} Un cadavre apoptotique clos ne libère rien dans une lumière
pleine de bactéries et ne provoque aucune inflammation~; une \omterm{def:b3:cell-cycle-apoptosis:apoptosis}{nécrose}
répandrait enzymes et signaux et enflammerait la muqueuse à chaque
cellule perdue.
\textbf{7.} $(40 - 8)/2 = \qty{16}{h}$.
\textbf{8.} $\log_{2}(40/8)\times \qty{10}{min} = 2.3\times 10 =
\qty{23}{min}$.
\textbf{9.} Période d'environ \qty{16.4}{h}~; Cdk1 active \qty{2.3}{\%}
du temps.
\textbf{10.} Ce n'est pas $k_{s}$ seul, mais le temps passé aux
seuils~: la cellule somatique attend au \omterm{def:b3:cell-cycle-apoptosis:checkpoints}{point de restriction} ses
facteurs de croissance et en G2/M l'état de son ADN, et la cellule
souche plus longtemps que la cellule de transit. L'œuf de grenouille
n'a pas de \omterm{def:b3:cell-cycle-apoptosis:checkpoints}{points de contrôle} et un cytoplasme approvisionné en tout,
de sorte que seul l'oscillateur nu fixe sa période.
\textbf{11.} Sur la branche basse, à un niveau de \omterm{def:b3:cell-cycle-apoptosis:cdk}{cycline} supérieur au
$C_{2}$ normal~: le \omterm{def:b3:cell-cycle-apoptosis:checkpoints}{point de contrôle} a décalé la courbe en S vers la
droite en inhibant Cdc25, de sorte que le saut n'a pas lieu. À la
réparation, Cdc25 est libérée, le seuil retombe sous la \omterm{def:b3:cell-cycle-apoptosis:cdk}{cycline}
accumulée, et la cellule entre aussitôt en mitose --- ce qui explique
que les cellules relâchées d'un \omterm{def:b3:cell-cycle-apoptosis:checkpoints}{point de contrôle} entrent en mitose de
façon synchrone.
\textbf{12.} Ni la \omterm{def:b3:cell-cycle-apoptosis:cdk}{cycline}~B ni la sécurine ne sont dégradées~: la
cellule reste en métaphase, cohésine intacte et Cdk1 active,
indéfiniment~; elle y meurt ou finit par en sortir avec un génome non
ségrégé.
\textbf{13.} $10^{11}$ divisions par jour.
\textbf{14.} Avec le \omterm{def:b3:cell-cycle-apoptosis:checkpoints}{point de contrôle} $10^{11}/5\times 10^{4} = 2\times 10^{6}$
filles aneuploïdes par jour~; sans lui, $2\times 10^{9}$.
\textbf{15.} $2\times 10^{4}$ par jour~; sur $80$ ans, $6\times 10^{8}$.
\textbf{16.} $10^{9}/50 = 2\times 10^{7}$ mauvaises ségrégations par
jour~: chaque jour la tumeur essaie vingt millions de caryotypes
nouveaux, et la sélection garde ceux qui croissent plus vite ou
résistent à un médicament.
\textbf{17.} Si chaque division ségrège mal, aucune lignée cellulaire
de l'embryon ne garde un génome euploïde et l'embryon meurt. Une perte
partielle donne une aneuploïdie occasionnelle dans des cellules qui y
survivent (surtout sans p53), instabilité chromosomique qui fournit de
la variation à une tumeur sans tuer l'organisme.
\textbf{18.} Une erreur méiotique met le chromosome surnuméraire dans
le gamète, de sorte que le zygote et toutes les cellules qui en
dérivent sont trisomiques~; une erreur mitotique survient dans une
cellule somatique et n'est héritée que par son clone.
\textbf{19.} $10^{11}\times \qty{1}{ng} = \qty{100}{g}$ par jour,
\qty{36}{kg} par an --- environ la moitié de la masse du corps chaque
année.
\textbf{20.} \qty{20}{g} de protéines par jour, soit environ le tiers de
l'apport alimentaire et quelques pour cent du renouvellement protéique
total du corps.
\textbf{21.} $10^{11}$ cadavres à $24$ par macrophage et par jour~:
$4\times
10^{9}$ macrophages à plein temps, soit \qty{4}{\%} du temps des
$10^{11}$ macrophages.
\textbf{22.} Un cadavre lysé libère des protéases, de l'ADN, de l'ATP
et d'autres signaux d'alarme qui causent une inflammation et peuvent
provoquer une auto-immunité contre les antigènes nucléaires~; le
cadavre intact s'annonce par la \omterm{def:b3:cell-cycle-apoptosis:apoptosis}{phosphatidylsérine} de son feuillet
externe (et par la perte de ses signaux «~ne me mangez pas~»), et
attire les phagocytes par des nucléotides libérés.
\textbf{23.} La \omterm{def:b3:cell-cycle-apoptosis:pathways}{voie intrinsèque} est bloquée à la mitochondrie. La \omterm{def:b3:cell-cycle-apoptosis:pathways}{voie
extrinsèque} fonctionne encore là où la caspase-8 active directement la
caspase-3, même si son amplification par Bid est perdue. Les cellules
s'accumulent faute de mourir plutôt qu'en se divisant plus vite ---
croissance lente --- jusqu'à ce qu'une lésion ultérieure (souvent
\emph{MYC}) y ajoute de la prolifération.
\textbf{24.} L'\omterm{def:b3:cell-cycle-apoptosis:apoptosis}{apoptose} prend des heures, exige de l'ATP et passe par
des étapes que l'on peut bloquer --- inhibiteurs de \omterm{def:b3:cell-cycle-apoptosis:apoptosis}{caspases}, seuil
mitochondrial --- de sorte que les neurones de la pénombre peuvent être
sauvés si l'on traite à temps~; ceux du cœur de la lésion meurent d'une
panne d'énergie en quelques minutes, sans aucun programme à
interrompre.
\textbf{25.} Période de l'oscillateur nu d'environ \qty{16}{h}~; quelque $2\times
10^{4}$ cellules aneuploïdes en division par jour dans un corps
normal~; \qty{100}{g} de cellules recyclées par \omterm{def:b3:cell-cycle-apoptosis:apoptosis}{apoptose} chaque jour.
\end{solution}
